\IEEEPARstart{T}{his} laboratory examines how a microcontroller acquires
physical signals and controls actuators. Systems 1--4 combine potentiometers,
light, motion, gas, distance, temperature and humidity sensors with visual or
audible responses. Systems 5--8 study a DC motor in simulation, LED dimming by
PWM, relay switching and moisture-controlled irrigation. The objective is to
connect each input and output correctly, convert raw readings to useful units,
implement the required thresholds and priorities, and compare the behavior
shown in simulations and physical checkpoints with the laboratory guides.
Arduino UNO is used for the specified Tinkercad simulations and MEGA2560 for
physical implementations. System 5's direct motor-to-pin circuit is strictly
a simulation exercise and is not a physical wiring recommendation.

\subsection{Theoretical Framework}

\subsubsection{Analog acquisition and scaling}
The MEGA2560 provides analog inputs and PWM-capable digital outputs
\cite{arduinoMega}. Under the 10-bit ADC and a 5-V reference specified in the
laboratory guide, a raw reading $n$ lies between 0 and 1023; the nominal
input voltage is $V\simeq 5n/1023$ V and the corresponding percentage is
$100n/1023$. A potentiometer connected as a voltage divider varies the
voltage at its center terminal. An LDR in a divider similarly converts a
change in illumination into a voltage. Its orientation in the divider must
be established before interpreting whether higher ADC counts mean more or
less light. For the four-LDR monitor, only valid channels contribute to the
mean, and a zero valid-channel count is handled as a fault rather than divided
by zero.

\subsubsection{Event sensing and priority}
A PIR sensor signals motion but does not itself identify a person; the
occupancy counter therefore uses a five-second lockout and a separate exit
button as required by System 3. An MQ2 analog reading is compared with an
established alarm threshold, with evacuation taking priority over normal
occupancy display. System 4 combines the HC-SR04 distance measurement with
temperature and humidity readings. The LM35 produces a nominal
$10\,\mathrm{mV}/{}^{\circ}\mathrm{C}$ analog output \cite{tiLM35}; the
DHT11 supplies digital environmental readings. Alternating LCD views should
be scheduled without unnecessarily clearing the display, while the alarm
temporarily overrides the second line.

\subsubsection{Actuation and electrical separation}
PWM controls the fraction of each cycle that a digital output is high. In
System 6 the 0--1023 potentiometer reading is mapped to the 0--255 input of
\texttt{analogWrite()}, and the LED requires a series current-limiting
resistor. A digital pin without PWM capability cannot provide gradual duty
cycle control. A motor draws more current than a microcontroller output is
designed to supply, so Systems 7 and 8 use an external supply and an
appropriate switching stage. The relay in System 7 separates the control
signal from the motor supply; the L9110S in System 8 drives the pump from
an external supply with a common ground for the logic signals. Irrigation
decisions depend on the measured dry and wet calibration endpoints, then on
the prescribed 20\% and 75\% thresholds and the 10-s/15-s actuation and
idle intervals.

\subsection{State of the Art}

Research on indoor occupancy sensing emphasizes that a PIR event is an
observation of motion, not a complete occupancy estimate. Emad-Ud-Din and
Wang review networked sensing methods and compare them by accuracy, failure
tolerance and observation rate \cite{emad2023}. System 3 is a simpler local
counter with one PIR and a manual exit input; its occupancy estimate can drift
if an entrance is missed or a button press does not match a real departure.
This motivates reporting the tested conditions and the limits of the count.

For irrigation, Garcia et al. survey sensor-based systems and identify soil
moisture as a common measurement used to make watering decisions
\cite{garcia2020}. A later systematic review by Abdelmoneim et al. discusses
the role of soil-moisture sensing and the risk of irrigating when sufficient
water is already present \cite{abdelmoneim2025}. System 8 demonstrates a
smaller, local threshold controller: it does not implement the network,
weather integration or crop-specific planning found in broader IoT systems.
Its evidence should therefore focus on sensor calibration, state transitions
and pump timing rather than claims of water savings that were not measured.
